\documentclass[10pt,a4paper]{amsart}
\usepackage[margin=19mm]{geometry}
\usepackage{amsmath,amssymb,mathtools}
\usepackage[scr=rsfs,cal=euler]{mathalfa}
\usepackage{xfrac}
\usepackage[unicode,hidelinks,bookmarks=false]{hyperref}
\newcommand{\ie}{i.e.\ }
\newcommand{\cf}{cf.\ }
\newcommand{\resp}{respectively\ }
\newcommand{\Subsection}{\subsection}
\providecommand{\nobreakdash}{\nobreak\hbox{-}}
\setlength{\emergencystretch}{2em}
\multlinegap0pt
\usepackage{mathtools}
\usepackage{enumerate}
\usepackage{amssymb}
\usepackage{tikz}
\newtheorem{lem}{Lemma}
\newtheorem{thm}{Theorem}
\def \Image{\operatorname{Image}}
\def\Trdim{\operatorname{Trdim}}
\def\her{\operatorname{her}}
\newcommand{\mc}{\mathcal}
\def\nuc{\operatorname{nuc}}
\title{On permanence of regularity properties II}
\author{Hyun Ho Lee}
\date{Source: arXiv:2603.07491v1 (CC BY 4.0)}
\begin{document}
\maketitle

\noindent\emph{Source and licence.} On permanence of regularity properties II. Version arXiv:2603.07491v1 (CC BY 4.0), distributed by arXiv under the Creative Commons Attribution 4.0 International licence, http://creativecommons.org/licenses/by/4.0/, as recorded in the arXiv licence metadata for this exact version and shown on the abstract page https://arxiv.org/abs/2603.07491v1. Full licence notice: \texttt{licenses/arxiv-2603.07491-CC-BY-4.0.txt}.\par\smallskip

\noindent\textit{Editorial scope.} Counted unit: the article's thm T:Trdim (source0.tex lines 423--425). The article's complete original proof of it is delivered with it (source0.tex lines 426--643, one \texttt{proof} environment of 218 lines). No mathematics is written, repaired or re-derived: the statement and the proof are the article's own lines, in the article's own notation, and no retained line is substituted or re-flowed.  Retained uncounted as the source-local context the proof needs: lines 353--360.  Nothing was omitted from any retained span, so the package carries no omission record.  No retained line contains a citation, so the wrapper carries no bibliography and no bibliography entry.  No retained line contains a figure environment, an image-inclusion command or drawing code, so no image is delivered and the package declares no asset dependency.  Editorial formatting only, in addition: an AMS \texttt{amsart} preamble that restates the article's own declarations and macros used by the retained lines, namely the environments \texttt{lem}, \texttt{thm} on separate counters, together with the standard packages those lines need.  The retained environments print in this document as Lemma 1, Theorem 1 in order of appearance, whereas the article prints the same items with the numbering of its own full document.  The complete native source of the article as distributed in the arXiv e-print for this exact version is bundled unchanged as \texttt{sources/195947/source0.tex}.
\begin{lem}(orthogonal lift of a c.p.c map)\label{L:orthogonallift}
Let $A$ be a separable nuclear unital $C\sp*$-algebra, $F$ a finite dimensional $C\sp*$-algebra and $\beta^{\omega}: F \to A^{\omega}$ a nonzero piecewise contractive  $m$-decomposable c.p. map. Put $D^{\omega}:=A^{\omega}\cap \beta^{\omega}(F)^{\perp}$. If $\Gamma^{\omega}: A \to D^{\omega}$ is c.p.c. then there exists a sequence of c.p.c. maps $\Gamma_n: A \to A$ such that 
\begin{enumerate}
\item $ [\Gamma_n(x)]=\Gamma^{\omega}(x)$ for all $x\in A$;
\item $ \Gamma_n(A) \subset A\cap \beta_n(F)^{\perp}$ for all $n$  
\end{enumerate}   
where $\beta^{\omega}=[\prod_{n=1}^{\infty} \beta_n]$ and $\beta_n: F \to A$ is a c.p. map  for each $n$. 
\end{lem}

\begin{thm}\label{T:Trdim}
Let $A$ be a nuclear simple unital infinite dimensional $C\sp*$-algebra and $B$ be a simple nuclear unital $C\sp*$-algebras and $\phi: A \to B $ a $^*$-homomorphism which is tracially sequentially-split by order zero. If $\Trdim_{\nuc} B \le m $, then $\Trdim_{\nuc} A \le m$. 
\end{thm}

\begin{proof} 
Given a finite set $\mc{F} \subset A$, $\epsilon>0$, $a_0 \in A_{+} \setminus \{ 0\}$, consider two nonzero positive elements $a_1, a_2 \in \her(a_0)$ such that $a_1 \perp a_2$ and $a_1 \sim a_2$.  consider a function $h$ defined by 
\[
 h(x)=\begin{cases}
 0 \quad & \text{if $0 \le x \le \eta$}\\
 \frac{x-\eta}{\eta} \quad & \text{if $\eta \le x \le 2\eta$}\\
 1 \quad &\text{if $x \ge 2 \eta$}
  \end{cases}
\] 
where $\eta >0$ is decided later.  Let $z:= h(a_1) \in A^{+}$. 
Then using the assumption we can take a c.p.c. order zero map $\psi: B \to A^{\omega}$  and $\psi(\mathbf{1_B})=g \in A^{\omega}\cap A'$ such that $\psi(\phi(a))=ag$,  $1- g \lesssim z \lesssim a_1$ since $z \in \her(a_1)$. 
We also consider $0< \delta < \epsilon$ which will be decided later.  Let $\mc{G}=\phi(\mc{F}) \subset B$, $\epsilon/6$, and $b_0=f(\phi(a_2))$ where $f\ge 0$ is a continuous function such that $f(0)=0, f(t)> 0$  for $t>0$. 
Since $\Trdim_{\nuc} B \le m$, there exist a finite dimensional $C\sp*$-algebra $F_0$, a c.p.c. map $\alpha_B: B \to F$, a nonzero piecewise $m$-decomposable c.p. map $\beta_B: F \to B$, a c.p.c. map $\gamma_B: B \to B\cap \beta(F)^{\perp}$, such that  
 \begin{enumerate}
\item $ \|y - (\gamma_B(y)+ \beta_B \circ \alpha_B (y) )\| < \epsilon/6$ for all $y\in \mc{G}$, and 
\item $\gamma_B(\mathbf{1_B}) \lesssim b_0$ in $A$. 
\end{enumerate}
Combining $\phi$ and $\psi$ with these maps, we define $\alpha: A \to F$ by $\alpha:= \alpha_B \circ \phi$,  $\beta^{\omega}: F \to A^{\omega}$ by $\beta^{\omega}:=\psi \circ \beta_B$, and $\gamma^{\omega}: A \to A^{\omega}$ by $\gamma^{\omega}:=\psi \circ \gamma_B \circ \phi$. Since $\psi$ is of order zero, it follows that \[\gamma^{\omega}(A) \perp \beta^{\omega}(F). \]

From (1), for any  $y \in \mc{G}$  
\begin{equation}
\| \psi(y) -(\psi(\gamma_B (y))+ \psi (\beta_B(\alpha_B(y)))) \| < \frac{\epsilon}{6}.
\end{equation}
So, for  any $x \in \mc{F}$
\begin{equation}
\| \psi(\phi(x)) -(\psi(\gamma_B (\phi(x)))+ \psi (\beta_B(\alpha_B(\phi(x))))) \| < \frac{\epsilon}{6}
\end{equation}
Thus, for  any $x \in \mc{F}$
\begin{equation}
\| xg -(\gamma^{\omega} (x)+  \beta^{\omega} (\alpha(x)) ) \| < \frac{\epsilon}{6}
\end{equation}
If we write $x=xg+ x(1-g)$, we have now the new remainder $\gamma^{\omega}(x) + x(1-g)$ replacing the role of $\gamma$ in the original definition. The problem is then the map $x \mapsto \gamma^{\omega}(x)+ x(1-g) \in A^{\omega}$ may not be orthogonal to $\beta^{\omega}(F)$. Let $b:=\beta^{\omega}(\mathbf{1_F}) \in A^{\omega}_{+} \setminus \{0\}$ and  $e=s(b)$ the support of $b$ in $(A^{\omega})^{**}$. It is well known that $\her(b)= A^{\omega}\cap e(A^{\omega})^{**}e$    and $\her^{\perp}(b)= A^{\omega}\cap (1-e)(A^{\omega})^{**}(1-e)= \overline{\{x \in A^{\omega} \mid xb=bx=0 \}}^{\| \,\|}$ due to Akemann and Pedersen. The latter is also a hereditary $C\sp*$-subalgebra.  
We note that $b$ is in the range of $\psi$, which means that $b \in \her(\psi(\mathbf{1_B}))=\her(g)$, and thus  $b=gc^{**}g$ for some $c^{**} \in (A^{\omega})^{**}$.  It follows that $b(1-g)=(1-g)b=0$. In fact, for every $x\in A$ $x(1-g)=(1-g)x \in \her^{\perp}(b)$ and in particular $1-g \in \her^{\perp}(b)$.  If we let $f_{\eta}(t)=\max\{0, t-\eta\}$, then $q:= f_{\eta}(1-g) \in \her^{\perp}(b)$ and $\| (1-g) -q\| \le \eta$. If we define $R^{\omega}(x):=q^{1/2}xq^{1/2}=qx$ for $x\in A$ and set $\eta= \frac{\epsilon}{ 6 \max\{ \|x\| \mid x\in \mc{F}\}}$, we can easily deduce that 
\begin{equation}
\| x - (\gamma^{\omega}(x) + R^{\omega}(x) + (\beta^{\omega}\circ \alpha)(c)) \| < \frac{\epsilon}{3}
\end{equation}
for all $x\in \mc{F}$. 
Moreover,  it is not hard to show that $R^{\omega}(x)\in \her^{\perp}(b)$ and $\Image( \beta^{\omega} )\subset \her(b)$. Therefore $R^{\omega}(A) \perp \beta^{\omega}(F)$.\\

On the other hand,  since $\psi$ is of order zero, 
\[
\gamma^{\omega}(\mathbf{1_A})=\psi(\gamma_B(\mathbf{1_B})) \lesssim \psi(f(\phi(a_2))) = f(a_2 g) \in \her(a_2). 
\]
Consequently, 
\begin{equation}\label{E:Cuntz1}
\gamma^{\omega}(\mathbf{1_A}) \lesssim a_2.
\end{equation}
Also, 
\begin{equation}\label{E:Cuntz2}
R^{\omega}(\mathbf{1_A})=q \le1-g \lesssim a_1
\end{equation}
So if we define $\Gamma^{\omega}:= \gamma^{\omega}+ R^{\omega}$, we obtain the following; 
\begin{equation}\label{E:condition(1)}
\|.x -(\Gamma^{\omega}(x)+ \beta^{\omega} \circ \alpha(x))  \| < \frac{\epsilon}{3}
\end{equation} 
\begin{equation}\label{E:condition(2)}
\Gamma^{\omega}(A) \perp \beta^{\omega}(F)
\end{equation}
\begin{equation}\label{E:condition(3)}
\Gamma^{\omega}(\mathbf{1_A}) \lesssim a_1 +a_2 \lesssim a_0  
\end{equation}
in $A^{\omega}$.\\  
Since $A$ nuclear, we can lift $\Gamma^{\omega}$ to $\Gamma: A \to \prod_{i=1}^{\infty} A$ and write $\Gamma= \prod_{n=1}^{\infty} \Gamma_n $ where $\Gamma_n: A \to A$ is a c.p. map by Choi-Effros. In addition, since $F$ is also nuclear, we can lift $\beta^{\omega}$ to $\widetilde{\beta}: F \to \prod_{n=1}^{\infty} A$ and write $\widetilde{\beta}=\prod_{n=1}^{\infty} \beta_n$ where $\beta_n: F \to A$ is a c.p.c map.    When we write $F=\oplus_{i=1}^m F_i $, the restriction map $\beta^{\omega}|_{F_i}$ is c.p.c order zero map for each $i$, we can write the c.p.c map $\beta_n=\sum_i \beta_{n, i}$ where $\beta_{n, i}: F_i \to A$ is also c.p.c order zero map. Combining (\ref{E:condition(2)}) with  Lemma \ref{L:orthogonallift}, we assume that $\Gamma_n(A) \perp \beta_n(F)$ for all $n$.\\  
\noindent
 From (\ref{E:condition(3)}), for each $k \in \mathbb{N}$, there exist $r^{(k)}=[(r^k_n)_{n=1}^{\infty}] \in A^{\omega}$ such that 
 \[  \| (r^{(k)})^* a_0 r^{(k)} - \Gamma(\mathbf{1_A}) \|  < \frac{1}{2^k} \quad \text{ in. $A^{\omega}$}. \] 
 Equivalently, 
 \[ \lim_{n \to \omega} \| (r^k_n)^* a_0 r^k_n - \Gamma_n (\mathbf{1_A}) \|_A < \frac{1}{2^k} \] 
 Therefore, $S_k:= \{j \in \mathbb{N}\mid \| (r^k_j)^* a_0  r^k_j - \Gamma_j (\mathbf{1_A})\|_A < \frac{1}{2^k}  \} \in \omega$. Note that $S_k$ has infinitely many elements for each $k$. \\
\noindent
For each $x \in \mc{F}$, from (\ref{E:condition(1)}), 
\[ \lim_{n\to \omega} \| x -(\Gamma_n(x)+ \beta_n \circ \alpha (x) )  \| < \frac{\epsilon}{3}. \]
Hence the set $U_x=\{n \mid \| x -(\Gamma_n(x)+ \beta_n \circ \alpha (x) )  \| < \frac{\epsilon}{3} \} \in \omega$. 
 Let
\[
U:=\bigcap_{x\in\mc{F}} U_x \in\omega.
\]
\noindent
For an arbitrary $\sigma >0$ (to be chosen later) and choose $k\in\mathbb N$ such that $2^{-k}<\sigma$. Since $U\in\omega$ as well, we have $U\cap S_k\in\omega$,
so we may pick $n\in U\cap S_k$. Then
\[
\|(r^{(k)}_{n})^*a_0 r^{(k)}_{n}-\Gamma_{n}(\mathbf{1_A})\|<2^{-k}<\sigma.
\]

Since $(r^{(k)}_{n})^*a_0 r^{(k)}_{n}\precsim a_0$ in $A$, the standard
perturbation-to-Cuntz-sub\-equivalence estimate yields
\[
(\Gamma_{n}(\mathbf{1_A})-\delta_1)_+\ \precsim\ (r^k_{n})^*a_0 r^k_{n}
\ \precsim\ a_0.
\]
Moreover, 
\begin{equation}\label{E:normestimate}
\| x -(\Gamma_{n}(x)+ \beta_{n} \circ \alpha (x) )  \| < \frac{\epsilon}{3}. 
\end{equation}
Now we set $\beta=\beta_n$, $\Gamma=\Gamma_n$ for our choice $n$. 

Now we need to deform $\Gamma$ to $\gamma$ such that $\gamma(\mathbf{1_A})= (\Gamma(\mathbf{1_A})-\sigma)_{+}$ and $\| \gamma(x) -\Gamma(x)\| < \dfrac{\epsilon}{3}$ for all $x \in \mc{F}$.  The following method might be known to expects;   
 First,  let 
\[H:=\her(h) \subset D_n=A\cap \beta_n(F)^{\perp} \] where $\Gamma(\mathbf{1_A})=h$.  Let $k_\sigma:=(f_\sigma(h))^{1/2} \in H$, where
\[
f_\sigma(t)=
\begin{cases}
0,& t=0,\\
\dfrac{(t-\sigma)_+}{t},& t > 0.
\end{cases}
\]
We define
\[
 \gamma(x):=k_{\sigma}\,\Gamma(x)\,k_{\sigma} \quad (x\in A).
\]

(1) Since $\gamma$ is obtained from $\Gamma$ by two-sided multiplication with a contraction,
it is completely positive and contractive and $k_{\sigma} \in H \subset D_n$ implies that $\gamma(A) \subset D_n$. 

\smallskip
(2) Since $k_\sigma$ commutes with $h$, we have
\[
\gamma(1_A)=k_\sigma\,h\,k_\sigma=g_\sigma(h)\,h.
\]
For  $t\ge 0$ one has $g_\sigma(t)\,t=(t-\sigma)_+$, hence by functional calculus
\[
\gamma(1_A)=(h-\sigma)_+.
\]

\smallskip
(3) Fix $x\in A$.
Let $\Gamma(x)=V^*\pi(x)V$ be a Stinespring dilation, so that $h=\Gamma(1_A)=V^*V$.
For $\varepsilon>0$ define the bounded operator
\[
W_\varepsilon:=V\,(h+\varepsilon)^{-1/2}.
\]
Then
\[
W_\varepsilon^*W_\varepsilon
=(h+\varepsilon)^{-1/2}V^*V(h+\varepsilon)^{-1/2}
=(h+\varepsilon)^{-1/2}h(h+\varepsilon)^{-1/2}
\le 1,
\]
so $\|W_\varepsilon\|\le 1$.
Set
\[
y_{x,\varepsilon}:=W_\varepsilon^*\pi(x)W_\varepsilon,
\]
which satisfies $\|y_{x,\varepsilon}\|\le \|x\|$.

Define
\[
s_\varepsilon:=h^{1/2}(h+\varepsilon)^{-1/2}\in C^*(h),\qquad 0\le s_\varepsilon\le 1.
\]
A direct computation gives
\[
s_\varepsilon\,\Gamma(x)\,s_\varepsilon
= h^{1/2}y_{x,\varepsilon}h^{1/2}.
\]
Moreover, $s_\varepsilon\to 1$ strongly on $\overline{\mathrm{Ran}(h)}$, and hence
\[
\|s_\varepsilon\Gamma(x)s_\varepsilon-\Gamma(x)\|\xrightarrow[\varepsilon\downarrow 0]{}0.
\]

Since $k_\sigma$ commutes with $h$ and
\[
k_\sigma h^{1/2}=(h-\sigma)_+^{1/2}
\]
by functional calculus, we have
\[
\gamma(x)
=k_\sigma\,h^{1/2}y_{x,\varepsilon}h^{1/2}k_\sigma
=(h-\sigma)_+^{1/2}y_{x,\varepsilon}(h-\sigma)_+^{1/2}.
\]

Let $d:=h^{1/2}-(h-\sigma)_+^{1/2}$. Then
\[
h^{1/2}y_{x,\varepsilon}h^{1/2}
-(h-\sigma)_+^{1/2}y_{x,\varepsilon}(h-\sigma)_+^{1/2}
= d\,y_{x,\varepsilon}\,h^{1/2}
+(h-\sigma)_+^{1/2}y_{x,\varepsilon}\,d.
\]
Taking norms and using $\|h^{1/2}\|\le 1$ and $\|(h-\sigma)_+^{1/2}\|\le 1$,
\[
\|h^{1/2}y_{x,\varepsilon}h^{1/2}
-(h-\sigma)_+^{1/2}y_{x,\varepsilon}(h-\sigma)_+^{1/2}\|
\le 2\|d\|\,\|y_{x,\varepsilon}\|.
\]
For  $t\in[0,1]$ one has
\[
0\le \sqrt{t}-\sqrt{(t-\sigma)_+}\le \sqrt{\sigma},
\]
hence $\|d\|\le \sqrt{\sigma}$ by functional calculus.
Since $\|y_{x,\varepsilon}\|\le \|x\|$, we obtain
\[
\|h^{1/2}y_{x,\varepsilon}h^{1/2}
-(h-\sigma)_+^{1/2}y_{x,\varepsilon}(h-\sigma)_+^{1/2}\|
\le 2\sqrt{\sigma}\,\|x\|.
\]

Letting $\varepsilon\downarrow 0$ yields
\[
\|\Gamma(x)-\gamma(x)\|\le 2\sqrt{\sigma}\,\|x\|,
\]
as required.
Now if $M:=\max\{ \|x\| \mid x \in \mc{F}\}> 0$, and we choose $0< \sigma < \frac{\varepsilon^2}{144M^2}$, then

\begin{equation}\label{E:perturbation}
\|\Gamma(x)-\gamma(x)\|<\frac{\varepsilon}{6}
\end{equation}
for all $x \in \mc{F}$. 
Combining two estimates (\ref{E:normestimate}) and (\ref{E:perturbation}) we obtain 
\[ 
\| x - (\gamma(x) + \beta \circ \alpha (x))\| < \frac{\epsilon}{2}< \epsilon
\]
for all $x \in \mc{F}$. 
Thus for given a finite set $\mc{F}$, $\varepsilon>0$, a nonzero positive element $a_0 \in A$, we have shown that there exist a finite dimensional $C\sp*$-algebra $F$, a c.p.c. map $\alpha: A \to F$, a nonzero piecewise contractive $m$-decomposable c.p. map $\beta: F \to A$, and a c.p.c map $\gamma: A \to A\cap \beta(F)^{\perp}$, such that 
\begin{enumerate}
\item $ \|x - (\gamma(x)+ \beta \circ \alpha (x)) \| < \epsilon$ for all $x\in \mc{F}$, and 
\item $\gamma(\mathbf{1_A}) \lesssim a_0$ in $A$. 
\end{enumerate}
Hence $\Trdim_{\nuc}A \le m$. 
\end{proof}

\end{document}
